\section{Exercices d'application}

\rubrique{Plan d'étude et domaine de définition}

\exo{}\diff{1}
Déterminer le domaine de définition des fonctions suivantes :
\[
f(x)=\frac{2x+1}{x^2-4},\qquad g(x)=\sqrt{3x-6},\qquad h(x)=\frac{\sqrt{x+1}}{x-3}.
\]

\exo{}\diff{1}
On considère la fonction \(f\) définie par \(f(x)=\frac{x^2-1}{x+2}\).
\begin{enumerate}[label=\alph*)]
\item Déterminer son domaine de définition.
\item Étudier la parité de \(f\).
\item Calculer les limites aux bornes du domaine.
\end{enumerate}

\rubrique{Limites et asymptotes}

\exo{}\diff{2}
Soit \(f(x)=\frac{2x^2-3x+1}{x-1}\).
\begin{enumerate}[label=\alph*)]
\item Déterminer le domaine de définition.
\item Calculer les limites aux bornes du domaine.
\item En déduire les asymptotes éventuelles.
\end{enumerate}

\exo{}\diff{2}
On donne \(f(x)=\frac{x^2+2x-3}{x+1}\).
\begin{enumerate}[label=\alph*)]
\item Déterminer les réels \(a,b,c\) tels que \(f(x)=ax+b+\frac{c}{x+1}\).
\item En déduire l'asymptote oblique à la courbe de \(f\).
\end{enumerate}

\rubrique{Tableau de variations et points particuliers}

\exo{}\diff{2}
Soit \(f(x)=x^3-3x+1\).
\begin{enumerate}[label=\alph*)]
\item Déterminer le domaine de définition.
\item Calculer la dérivée \(f'(x)\) et étudier son signe.
\item Dresser le tableau de variations de \(f\).
\item Déterminer les extremums locaux.
\end{enumerate}

\exo{}\diff{2}
Étudier les variations de \(f(x)=\frac{x^2+1}{x-2}\) sur son domaine de définition.

\rubrique{Étude complète de fonctions polynômes}

\exo{}\diff{2}
Soit \(f(x)=x^3-6x^2+9x-4\).
\begin{enumerate}[label=\alph*)]
\item Déterminer le domaine de définition.
\item Calculer la dérivée et étudier son signe.
\item Dresser le tableau de variations.
\item Déterminer les points d'intersection avec les axes.
\item Tracer la courbe représentative de \(f\).
\end{enumerate}

\exo{}\diff{3}
On considère \(f(x)=x^4-4x^2+3\).
\begin{enumerate}[label=\alph*)]
\item Étudier la parité de \(f\).
\item Déterminer les limites en \(+\infty\) et \(-\infty\).
\item Calculer la dérivée et étudier son signe.
\item Dresser le tableau de variations.
\item Tracer la courbe.
\end{enumerate}

\rubrique{Étude complète de fonctions rationnelles}

\exo{}\diff{2}
Soit \(f(x)=\frac{x^2-4}{x-1}\).
\begin{enumerate}[label=\alph*)]
\item Déterminer le domaine de définition.
\item Calculer les limites aux bornes du domaine.
\item Déterminer les asymptotes.
\item Calculer la dérivée et étudier son signe.
\item Dresser le tableau de variations.
\item Tracer la courbe.
\end{enumerate}

\exo{}\diff{3}
Étudier complètement la fonction \(f(x)=\frac{x^2+3x+2}{x+2}\) (domaine, limites, asymptotes, variations, tracé).

\rubrique{Étude complète de fonctions irrationnelles simples}

\exo{}\diff{2}
Soit \(f(x)=\sqrt{x^2-4}\).
\begin{enumerate}[label=\alph*)]
\item Déterminer le domaine de définition.
\item Étudier la parité.
\item Calculer les limites en \(+\infty\) et \(-\infty\).
\item Déterminer les asymptotes éventuelles.
\item Calculer la dérivée et étudier son signe.
\item Dresser le tableau de variations.
\item Tracer la courbe.
\end{enumerate}

\exo{}\diff{3}
On considère \(f(x)=\sqrt{x+1}-\sqrt{x-1}\).
\begin{enumerate}[label=\alph*)]
\item Déterminer le domaine de définition.
\item Calculer les limites aux bornes du domaine.
\item Étudier les variations de \(f\).
\item Tracer la courbe.
\end{enumerate}

\section{Exercices d'entraînement}

\rubrique{Domaine de définition et parité}

\exo{}\diff{1}
Déterminer le domaine de définition de :
\[
f(x)=\frac{3x-5}{x^2+2x-3},\quad g(x)=\sqrt{5-2x},\quad h(x)=\frac{\sqrt{x-2}}{x^2-9}.
\]

\exo{}\diff{1}
Étudier la parité des fonctions suivantes :
\[
f(x)=x^4-3x^2+2,\quad g(x)=x^3-2x,\quad h(x)=\frac{x}{x^2+1}.
\]

\exo{}\diff{2}
Soit \(f(x)=\frac{x^2+1}{x^2-1}\). Étudier sa parité et déterminer son domaine de définition.

\rubrique{Limites et asymptotes}

\exo{}\diff{2}
Calculer les limites suivantes :
\[
\lim_{x\to 2}\frac{x^2-4}{x-2},\quad \lim_{x\to +\infty}\frac{3x^2-2x+1}{x^2+5},\quad \lim_{x\to -\infty}\frac{2x^3-5x}{x^2+1}.
\]

\exo{}\diff{2}
Déterminer les asymptotes de \(f(x)=\frac{2x^2+3x-5}{x+1}\).

\exo{}\diff{2}
Soit \(f(x)=\frac{x^2-2x+3}{x-2}\). Déterminer les réels \(a,b,c\) tels que \(f(x)=ax+b+\frac{c}{x-2}\). En déduire l'asymptote oblique.

\exo{}\diff{3}
On considère \(f(x)=\frac{x^3-2x^2+1}{x^2-1}\). Déterminer les asymptotes.

\rubrique{Tableau de variations}

\exo{}\diff{2}
Étudier les variations de \(f(x)=x^3-3x^2+2\).

\exo{}\diff{2}
Soit \(f(x)=\frac{x^2-2x+1}{x-3}\). Calculer la dérivée et étudier son signe.

\exo{}\diff{2}
Dresser le tableau de variations de \(f(x)=\sqrt{x^2-2x+5}\).

\exo{}\diff{3}
Étudier les variations de \(f(x)=\frac{x^2+2x+2}{x+1}\).

\rubrique{Études complètes}

\exo{}\diff{2}
Étudier complètement \(f(x)=x^3-3x^2-9x+5\) (domaine, limites, variations, points particuliers, tracé).

\exo{}\diff{2}
Soit \(f(x)=\frac{x^2-3x+2}{x-1}\). Faire une étude complète.

\exo{}\diff{2}
Étudier \(f(x)=\sqrt{x^2-2x}\).

\exo{}\diff{3}
On donne \(f(x)=\frac{x^3-4x}{x^2-1}\). Étudier complètement cette fonction.

\exo{}\diff{3}
Étudier \(f(x)=\sqrt{x+2}-\sqrt{x-2}\).

\exo{}\diff{3}
Soit \(f(x)=\frac{x^2+2x+3}{x+1}\). Déterminer les positions relatives de la courbe et de son asymptote oblique.

\exo{}\diff{3}
Étudier complètement \(f(x)=\frac{x^2-2x+4}{x-2}\).

\exo{}\diff{3}
On considère \(f(x)=\sqrt{x^2+2x+3}\). Étudier les variations et tracer la courbe.

\exo{}\diff{3}
Étudier \(f(x)=\frac{x^3-3x+2}{x^2-1}\).

\section{Problèmes type Probatoire/Bac}

\sujetbac[Étude d'une fonction rationnelle]
On considère la fonction \(f\) définie par \(f(x)=\frac{x^2-2x+3}{x-1}\).
\begin{enumerate}[label=\arabic*)]
\item Déterminer le domaine de définition de \(f\).
\item Calculer les limites aux bornes du domaine. En déduire les asymptotes.
\item Déterminer les réels \(a,b,c\) tels que \(f(x)=ax+b+\frac{c}{x-1}\).
\item Calculer la dérivée \(f'(x)\) et étudier son signe.
\item Dresser le tableau de variations de \(f\).
\item Déterminer les points d'intersection de la courbe avec les axes.
\item Étudier la position relative de la courbe et de son asymptote oblique.
\item Tracer la courbe représentative de \(f\) dans un repère orthonormé.
\end{enumerate}

\sujetbac[Étude d'une fonction irrationnelle]
Soit \(f(x)=\sqrt{x^2-2x+2}\).
\begin{enumerate}[label=\arabic*)]
\item Déterminer le domaine de définition de \(f\).
\item Étudier la parité de \(f\).
\item Calculer les limites en \(+\infty\) et \(-\infty\).
\item Déterminer les asymptotes éventuelles.
\item Calculer la dérivée \(f'(x)\) et étudier son signe.
\item Dresser le tableau de variations.
\item Tracer la courbe représentative de \(f\).
\end{enumerate}

\sujetbac[Étude d'une fonction polynôme]
On considère \(f(x)=x^3-3x^2-9x+11\).
\begin{enumerate}[label=\arabic*)]
\item Déterminer le domaine de définition.
\item Calculer les limites en \(+\infty\) et \(-\infty\).
\item Calculer la dérivée \(f'(x)\) et étudier son signe.
\item Dresser le tableau de variations.
\item Déterminer les extremums locaux.
\item Déterminer les points d'intersection avec les axes.
\item Tracer la courbe.
\end{enumerate}

\sujetbac[Étude d'une fonction rationnelle avec asymptote oblique]
Soit \(f(x)=\frac{2x^2-3x+4}{x-1}\).
\begin{enumerate}[label=\arabic*)]
\item Déterminer le domaine de définition.
\item Calculer les limites aux bornes du domaine.
\item Déterminer les asymptotes.
\item Calculer la dérivée et étudier son signe.
\item Dresser le tableau de variations.
\item Tracer la courbe.
\end{enumerate}

\section{Corrigés}

\corrige{de l'exercice 1}
\begin{enumerate}[label=\alph*)]
\item \(f(x)=\frac{2x+1}{x^2-4}\). Le dénominateur s'annule pour \(x^2-4=0\iff x=2\) ou \(x=-2\). Donc \(\mathcal{D}_f=\R\setminus\{-2,2\}\).
\item \(g(x)=\sqrt{3x-6}\). Il faut \(3x-6\ge0\iff x\ge2\). Donc \(\mathcal{D}_g=[2,+\infty[\).
\item \(h(x)=\frac{\sqrt{x+1}}{x-3}\). Il faut \(x+1\ge0\iff x\ge-1\) et \(x-3\neq0\iff x\neq3\). Donc \(\mathcal{D}_h=[-1,3[\cup]3,+\infty[\).
\end{enumerate}

\corrige{de l'exercice 2}
\begin{enumerate}[label=\alph*)]
\item \(f(x)=\frac{x^2-1}{x+2}\). Le dénominateur s'annule pour \(x=-2\). Donc \(\mathcal{D}_f=\R\setminus\{-2\}\).
\item \(f(-x)=\frac{(-x)^2-1}{-x+2}=\frac{x^2-1}{-x+2}\). On a \(f(-x)\neq f(x)\) et \(f(-x)\neq -f(x)\). Donc \(f\) n'est ni paire ni impaire.
\item \(\lim_{x\to -\infty}f(x)=\lim_{x\to -\infty}\frac{x^2}{x}= -\infty\). \(\lim_{x\to +\infty}f(x)=+\infty\). \(\lim_{x\to -2^-}f(x)=+\infty\) (car numérateur tend vers 3, dénominateur tend vers $0^-$). \(\lim_{x\to -2^+}f(x)=-\infty\).
\end{enumerate}

\corrige{de l'exercice 3}
\begin{enumerate}[label=\alph*)]
\item \(f(x)=\frac{2x^2-3x+1}{x-1}\). Le dénominateur s'annule pour \(x=1\). Donc \(\mathcal{D}_f=\R\setminus\{1\}\).
\item \(\lim_{x\to -\infty}f(x)=\lim_{x\to -\infty}\frac{2x^2}{x}= -\infty\). \(\lim_{x\to +\infty}f(x)=+\infty\). \(\lim_{x\to 1^-}f(x)=-\infty\) (numérateur tend vers 0, dénominateur tend vers $0^-$). \(\lim_{x\to 1^+}f(x)=+\infty\).
\item Asymptote verticale : \(x=1\). Asymptote oblique : on effectue la division euclidienne : \(2x^2-3x+1=(x-1)(2x-1)+0\). Donc \(f(x)=2x-1\). La droite \(y=2x-1\) est asymptote oblique.
\end{enumerate}

\corrige{de l'exercice 4}
\begin{enumerate}[label=\alph*)]
\item On effectue la division : \(x^2+2x-3=(x+1)(x+1)-4\). Donc \(f(x)=x+1-\frac{4}{x+1}\). Ainsi \(a=1\), \(b=1\), \(c=-4\).
\item L'asymptote oblique est la droite \(y=x+1\).
\end{enumerate}

\corrige{de l'exercice 5}
\begin{enumerate}[label=\alph*)]
\item \(\mathcal{D}_f=\R\).
\item \(f'(x)=3x^2-3=3(x^2-1)=3(x-1)(x+1)\). Le signe de \(f'(x)\) est positif pour \(x<-1\) ou \(x>1\), négatif pour \(-1<x<1\).
\item Tableau de variations :
\[
\begin{array}{c|ccccccccccc}
x & -\infty & & -1 & & 1 & & +\infty \\
\hline
f'(x) & & + & 0 & - & 0 & + & \\
\hline
f(x) & -\infty & \nearrow & 3 & \searrow & -1 & \nearrow & +\infty
\end{array}
\]
\item Maximum local : \(f(-1)=3\). Minimum local : \(f(1)=-1\).
\end{enumerate}

\corrige{de l'exercice 6}
\(f(x)=\frac{x^2+1}{x-2}\). \(\mathcal{D}_f=\R\setminus\{2\}\). \(f'(x)=\frac{2x(x-2)-(x^2+1)}{(x-2)^2}=\frac{2x^2-4x-x^2-1}{(x-2)^2}=\frac{x^2-4x-1}{(x-2)^2}\). Le numérateur s'annule pour \(x=2\pm\sqrt{5}\). Signe : positif à l'extérieur des racines, négatif entre. Tableau :
\[
\begin{array}{c|ccccccccccc}
x & -\infty & & 2-\sqrt{5} & & 2 & & 2+\sqrt{5} & & +\infty \\
\hline
f'(x) & & + & 0 & - & \text{non déf} & - & 0 & + & \\
\hline
f(x) & -\infty & \nearrow & \text{max} & \searrow & \text{asymp} & \searrow & \text{min} & \nearrow & +\infty
\end{array}
\]

\corrige{de l'exercice 7}
\begin{enumerate}[label=\alph*)]
\item \(\mathcal{D}_f=\R\).
\item \(f'(x)=3x^2-12x+9=3(x^2-4x+3)=3(x-1)(x-3)\). Signe : positif pour \(x<1\) ou \(x>3\), négatif pour \(1<x<3\).
\item Tableau :
\[
\begin{array}{c|ccccccccccc}
x & -\infty & & 1 & & 3 & & +\infty \\
\hline
f'(x) & & + & 0 & - & 0 & + & \\
\hline
f(x) & -\infty & \nearrow & 0 & \searrow & -4 & \nearrow & +\infty
\end{array}
\]
\item Intersection avec l'axe des ordonnées : \(f(0)=-4\). Avec l'axe des abscisses : \(f(x)=0\iff x^3-6x^2+9x-4=0\). On remarque \(x=1\) est racine : \((x-1)(x^2-5x+4)=0\iff (x-1)^2(x-4)=0\). Donc points \((1,0)\) et \((4,0)\).
\item Tracé : courbe cubique passant par ces points.
\end{enumerate}

\corrige{de l'exercice 8}
\begin{enumerate}[label=\alph*)]
\item \(f(-x)=(-x)^4-4(-x)^2+3=x^4-4x^2+3=f(x)\). Donc \(f\) est paire.
\item \(\lim_{x\to +\infty}f(x)=+\infty\), \(\lim_{x\to -\infty}f(x)=+\infty\).
\item \(f'(x)=4x^3-8x=4x(x^2-2)=4x(x-\sqrt{2})(x+\sqrt{2})\). Signe : \(f'(x)>0\) pour \(-\sqrt{2}<x<0\) ou \(x>\sqrt{2}\) ; \(f'(x)<0\) pour \(x<-\sqrt{2}\) ou \(0<x<\sqrt{2}\).
\item Tableau :
\[
\begin{array}{c|ccccccccccc}
x & -\infty & & -\sqrt{2} & & 0 & & \sqrt{2} & & +\infty \\
\hline
f'(x) & & - & 0 & + & 0 & - & 0 & + & \\
\hline
f(x) & +\infty & \searrow & -1 & \nearrow & 3 & \searrow & -1 & \nearrow & +\infty
\end{array}
\]
\item Tracé : courbe symétrique par rapport à l'axe des ordonnées.
\end{enumerate}

\corrige{de l'exercice 9}
\begin{enumerate}[label=\alph*)]
\item \(\mathcal{D}_f=\R\setminus\{1\}\).
\item \(\lim_{x\to -\infty}f(x)=-\infty\), \(\lim_{x\to +\infty}f(x)=+\infty\). \(\lim_{x\to 1^-}f(x)=+\infty\), \(\lim_{x\to 1^+}f(x)=-\infty\).
\item Asymptote verticale : \(x=1\). Asymptote oblique : division euclidienne : \(x^2-4=(x-1)(x+1)-3\). Donc \(f(x)=x+1-\frac{3}{x-1}\). Asymptote oblique : \(y=x+1\).
\item \(f'(x)=1+\frac{3}{(x-1)^2}>0\) pour tout \(x\neq1\).
\item Tableau :
\[
\begin{array}{c|ccccccccccc}
x & -\infty & & 1 & & +\infty \\
\hline
f'(x) & & + & \text{non déf} & + & \\
\hline
f(x) & -\infty & \nearrow & \text{asymp} & \nearrow & +\infty
\end{array}
\]
\item Tracé : hyperbole avec asymptotes.
\end{enumerate}

\corrige{de l'exercice 10}
\(f(x)=\frac{x^2+3x+2}{x+2}\). \(\mathcal{D}_f=\R\setminus\{-2\}\). Factorisation : \(x^2+3x+2=(x+1)(x+2)\). Donc \(f(x)=x+1\) pour \(x\neq-2\). La courbe est la droite \(y=x+1\) privée du point \((-2,-1)\). Asymptote verticale : aucune (car simplification). Limites : \(\lim_{x\to -2}f(x)=-1\). Variations : \(f'(x)=1>0\). Tableau : \(f\) croissante sur chaque intervalle.

\corrige{de l'exercice 11}
\begin{enumerate}[label=\alph*)]
\item Il faut \(x^2-4\ge0\iff x\le-2\) ou \(x\ge2\). Donc \(\mathcal{D}_f=]-\infty,-2]\cup[2,+\infty[\).
\item \(f(-x)=\sqrt{(-x)^2-4}=\sqrt{x^2-4}=f(x)\). Donc \(f\) est paire.
\item \(\lim_{x\to +\infty}f(x)=+\infty\), \(\lim_{x\to -\infty}f(x)=+\infty\).
\item \(\lim_{x\to +\infty}\frac{f(x)}{x}=\lim_{x\to +\infty}\frac{\sqrt{x^2-4}}{x}=1\). \(\lim_{x\to +\infty}(f(x)-x)=\lim_{x\to +\infty}(\sqrt{x^2-4}-x)=0\). Donc asymptote oblique \(y=x\) en \(+\infty\). Par symétrie, \(y=-x\) en \(-\infty\).
\item \(f'(x)=\frac{x}{\sqrt{x^2-4}}\). Signe : positif pour \(x>2\), négatif pour \(x<-2\).
\item Tableau :
\[
\begin{array}{c|ccccccccccc}
x & -\infty & & -2 & & 2 & & +\infty \\
\hline
f'(x) & & - & \text{non déf} & & \text{non déf} & + & \\
\hline
f(x) & +\infty & \searrow & 0 & & 0 & \nearrow & +\infty
\end{array}
\]
\item Tracé : deux branches symétriques.
\end{enumerate}

\corrige{de l'exercice 12}
\begin{enumerate}[label=\alph*)]
\item Il faut \(x+1\ge0\) et \(x-1\ge0\iff x\ge1\). Donc \(\mathcal{D}_f=[1,+\infty[\).
\item \(\lim_{x\to 1^+}f(x)=\sqrt{2}-0=\sqrt{2}\). \(\lim_{x\to +\infty}f(x)=+\infty\).
\item \(f'(x)=\frac{1}{2\sqrt{x+1}}-\frac{1}{2\sqrt{x-1}}<0\) pour \(x>1\). Donc \(f\) est strictement décroissante.
\item Tracé : courbe décroissante de \(\sqrt{2}\) à \(+\infty\).
\end{enumerate}

\corrige{du problème 1}
\begin{enumerate}[label=\arabic*)]
\item \(\mathcal{D}_f=\R\setminus\{1\}\).
\item \(\lim_{x\to -\infty}f(x)=-\infty\), \(\lim_{x\to +\infty}f(x)=+\infty\). \(\lim_{x\to 1^-}f(x)=-\infty\), \(\lim_{x\to 1^+}f(x)=+\infty\). Asymptote verticale : \(x=1\).
\item Division : \(x^2-2x+3=(x-1)(x-1)+2\). Donc \(f(x)=x-1+\frac{2}{x-1}\).
\item \(f'(x)=1-\frac{2}{(x-1)^2}=\frac{(x-1)^2-2}{(x-1)^2}=\frac{x^2-2x-1}{(x-1)^2}\). Racines : \(x=1\pm\sqrt{2}\). Signe : positif pour \(x<1-\sqrt{2}\) ou \(x>1+\sqrt{2}\), négatif entre.
\item Tableau :
\[
\begin{array}{c|ccccccccccc}
x & -\infty & & 1-\sqrt{2} & & 1 & & 1+\sqrt{2} & & +\infty \\
\hline
f'(x) & & + & 0 & - & \text{non déf} & - & 0 & + & \\
\hline
f(x) & -\infty & \nearrow & \text{max} & \searrow & \text{asymp} & \searrow & \text{min} & \nearrow & +\infty
\end{array}
\]
\item Intersection avec l'axe des ordonnées : \(f(0)=-3\). Avec l'axe des abscisses : \(f(x)=0\iff x^2-2x+3=0\), pas de solution (discriminant négatif). Donc pas d'intersection.
\item Position relative : \(f(x)-(x-1)=\frac{2}{x-1}\). Pour \(x>1\), \(\frac{2}{x-1}>0\) donc courbe au-dessus de l'asymptote. Pour \(x<1\), courbe en dessous.
\item Tracé : hyperbole avec asymptotes.
\end{enumerate}

\corrige{du problème 2}
\begin{enumerate}[label=\arabic*)]
\item Il faut \(x^2-2x+2\ge0\). Discriminant \(\Delta=4-8=-4<0\), donc toujours positif. \(\mathcal{D}_f=\R\).
\item \(f(-x)=\sqrt{x^2+2x+2}\neq f(x)\) et \(\neq -f(x)\). Donc ni paire ni impaire.
\item \(\lim_{x\to +\infty}f(x)=+\infty\), \(\lim_{x\to -\infty}f(x)=+\infty\).
\item \(\lim_{x\to +\infty}\frac{f(x)}{x}=1\), \(\lim_{x\to +\infty}(f(x)-x)=\lim_{x\to +\infty}\frac{-2x+2}{\sqrt{x^2-2x+2}+x}=-1\). Donc asymptote oblique \(y=x-1\) en \(+\infty\). En \(-\infty\), \(\lim_{x\to -\infty}\frac{f(x)}{x}=-1\), \(\lim_{x\to -\infty}(f(x)+x)=1\), donc asymptote \(y=-x+1\).
\item \(f'(x)=\frac{x-1}{\sqrt{x^2-2x+2}}\). Signe : négatif pour \(x<1\), positif pour \(x>1\).
\item Tableau :
\[
\begin{array}{c|ccccccccccc}
x & -\infty & & 1 & & +\infty \\
\hline
f'(x) & & - & 0 & + & \\
\hline
f(x) & +\infty & \searrow & 1 & \nearrow & +\infty
\end{array}
\]
\item Tracé : courbe avec minimum en \(x=1\).
\end{enumerate}

\corrige{du problème 3}
\begin{enumerate}[label=\arabic*)]
\item \(\mathcal{D}_f=\R\).
\item \(\lim_{x\to -\infty}f(x)=-\infty\), \(\lim_{x\to +\infty}f(x)=+\infty\).
\item \(f'(x)=3x^2-6x-9=3(x^2-2x-3)=3(x-3)(x+1)\). Signe : positif pour \(x<-1\) ou \(x>3\), négatif pour \(-1<x<3\).
\item Tableau :
\[
\begin{array}{c|ccccccccccc}
x & -\infty & & -1 & & 3 & & +\infty \\
\hline
f'(x) & & + & 0 & - & 0 & + & \\
\hline
f(x) & -\infty & \nearrow & 16 & \searrow & -16 & \nearrow & +\infty
\end{array}
\]
\item Maximum local : \(f(-1)=16\). Minimum local : \(f(3)=-16\).
\item Intersection avec l'axe des ordonnées : \(f(0)=11\). Avec l'axe des abscisses : \(f(x)=0\). On cherche des racines évidentes : \(x=1\) donne \(1-3-9+11=0\). Donc \((x-1)(x^2-2x-11)=0\). Racines : \(x=1\) et \(x=1\pm2\sqrt{3}\). Points : \((1,0)\), \((1+2\sqrt{3},0)\), \((1-2\sqrt{3},0)\).
\item Tracé : courbe cubique.
\end{enumerate}

\corrige{du problème 4}
\begin{enumerate}[label=\arabic*)]
\item \(\mathcal{D}_f=\R\setminus\{1\}\).
\item \(\lim_{x\to -\infty}f(x)=-\infty\), \(\lim_{x\to +\infty}f(x)=+\infty\). \(\lim_{x\to 1^-}f(x)=-\infty\), \(\lim_{x\to 1^+}f(x)=+\infty\).
\item Asymptote verticale : \(x=1\). Asymptote oblique : division : \(2x^2-3x+4=(x-1)(2x-1)+3\). Donc \(f(x)=2x-1+\frac{3}{x-1}\). Asymptote oblique : \(y=2x-1\).
\item \(f'(x)=2-\frac{3}{(x-1)^2}=\frac{2(x-1)^2-3}{(x-1)^2}=\frac{2x^2-4x-1}{(x-1)^2}\). Racines : \(x=\frac{2\pm\sqrt{6}}{2}\). Signe : positif pour \(x<\frac{2-\sqrt{6}}{2}\) ou \(x>\frac{2+\sqrt{6}}{2}\), négatif entre.
\item Tableau :
\[
\begin{array}{c|ccccccccccc}
x & -\infty & & \frac{2-\sqrt{6}}{2} & & 1 & & \frac{2+\sqrt{6}}{2} & & +\infty \\
\hline
f'(x) & & + & 0 & - & \text{non déf} & - & 0 & + & \\
\hline
f(x) & -\infty & \nearrow & \text{max} & \searrow & \text{asymp} & \searrow & \text{min} & \nearrow & +\infty
\end{array}
\]
\item Tracé : hyperbole avec asymptotes.
\end{enumerate}